%\documentclass[a4paper, twocolumn]{article}
\documentclass[11pt]{article}
\pagestyle{empty}

%\usepackage[a6paper]{geometry}
\usepackage[portuguese]{babel}
\usepackage[latin1]{inputenc}
\usepackage[T1]{fontenc}

\usepackage{amssymb,amsmath}
%\usepackage{euler}
\everymath{\displaystyle}

\DeclareMathOperator{\N}{\mathbb{N}}
\DeclareMathOperator{\Z}{\mathbb{Z}}
\DeclareMathOperator{\R}{\mathbb{R}}
\DeclareMathOperator{\arccot}{arccot}
\DeclareMathOperator{\arcsec}{arcsec}
\DeclareMathOperator{\arccsc}{arccsc}
\DeclareMathOperator{\negate}
    {\negthickspace\negthickspace\negthickspace\diagup\negthickspace}
\DeclareMathOperator{\dom}{dom}
\DeclareMathOperator{\im}{im}
\providecommand{\abs}[1]{\lvert#1\rvert}
\providecommand{\ddx}{\frac{d}{dx}}

\author{Andr� Wagner}
\title{C�lculo}
\date{}

\begin{document}

\section{Trigonometria}

\begin{enumerate}

  \item $ \sin \theta = \frac{O}{H} $
  \item $ \cos \theta = \frac{A}{H} $
  \item $ \tan \theta = \frac{O}{A} $
  \item $ \csc \theta = \frac{H}{O} $
  \item $ \sec \theta = \frac{H}{A} $
  \item $ \cot \theta = \frac{A}{O} $

\end{enumerate}

\section{Identidades trigonom�tricas}

\begin{enumerate}

  \item $ \csc\theta = \frac{1}{\sin\theta} $
  \item $ \sec\theta = \frac{1}{\cos\theta} $
  \item $ \cot\theta = \frac{1}{\tan\theta} $
  \item $ \tan\theta = \frac{\sin\theta}{\cos\theta} $
  \item $ \cot\theta = \frac{\cos\theta}{\sin\theta} $
  \item $ \sin^2\theta + \cos^2\theta = 1 $
  \item $ \tan^2\theta + 1 = \sec^2\theta $
  \item $ 1+ \cot^2\theta = \csc^2\theta $
  \item $ \sin^2\theta = \frac{1 - \cos 2\theta}{2} $
  \item $ \cos^2\theta = \frac{1 + \cos 2\theta}{2} $
  \item $ \sin 2\theta = 2\sin\theta\cos\theta $
  \item $ \cos 2\theta = 2\cos^2\theta - 1 $
  \item $ \sin(-\theta) = -\sin\theta $
  \item $ \cos(-\theta) = -\cos\theta $
  \item $ \tan(-\theta) = -\tan\theta $

\end{enumerate}

\section{Derivadas}

\begin{enumerate}

  \item $ \ddx(uv) = u'v + v'u $
  \item $ \ddx\left(\frac{u}{v}\right) = \frac{u'v - v'u}{v^2} $
  \item $ \ddx c = 0 $
  \item $ \ddx x = 1 $
  \item $ \ddx cx = c $
  \item $ \ddx x^n = nx^{n-1} $
  \item $ \ddx e^x = e^x $
  \item $ \ddx \ln x = \frac{1}{x} $
  \item $ \ddx a^x = a^x \ln a $
  \item $ \ddx \log_a x = \frac{1}{x} \cdot \frac{1}{\ln a} $
  \item $ \ddx \sin x = \cos x $
  \item $ \ddx \cos x = -\sin x $
  \item $ \ddx \tan x = \sec^2 x $
  \item $ \ddx \cot x = -\csc^2 x $
  \item $ \ddx \sec x = \sec x \tan x $
  \item $ \ddx \csc x = -\csc x \cot x $
  \item $ \ddx \arcsin x = \frac{1}{\sqrt{1-x^2}} $
  \item $ \ddx \arccos x = \frac{-1}{1+x^2} $
  \item $ \ddx \arctan x = \frac{1}{1+x^2} $
  \item $ \ddx \arccot x = \frac{-1}{1+x^2} $
  \item $ \ddx \arcsec x = \frac{1}{\abs{x}\sqrt{x^2-1}} $
  \item $ \ddx \arccsc x = \frac{-1}{\abs{x}\sqrt{x^2-1}} $

\end{enumerate}

\section{Integrais}

\begin{enumerate}

  \item $ \int dx = x + C $
  \item $ \int x^n dx = \frac{x^{n+1}}{n+1} + C (n \neq 1) $
  \item $ \i\mathbb{R}nt e^x dx = e^x + C $
  \item $ \int \frac{dx}{x} = \ln\abs{x} + C $
  \item $ \int a^x dx = \frac{a^x}{\ln a} + C $
  \item $ \int \ln x dx = x(\ln x-1) + C $
  \item $ \int \sin x dx = -\cos x + C $
  \item $ \int \cos x dx = \sin x + C $
  \item $ \int \tan x dx = -\ln\abs{\cos x} + C $
  \item $ \int \cot x dx = \ln\abs{\sin x} + C $
  \item $ \int \sec x dx = \ln\abs{\sec x + \tan x} + C $
  \item $ \int \csc x dx = -\ln\abs{\csc x + \cot x} + C $
  \item $ \int \sec^2 x dx = \tan x + C $
  \item $ \int \csc^2 x dx = -\cot x + C $
  \item $ \int \sec x \tan x dx = \sec x + C $
  \item $ \int \csc x \cot x dx = -\csc x + C $
  \item $ \int \frac{dx}{\sqrt{a^2-x^2}} = \arcsin\frac{x}{a} + C $
  \item $ \int \frac{dx}{x\sqrt{x^2-a^2}} = \frac{1}{a}\arcsec\frac{\abs{x}}{a} + C $
  \item $ \int \frac{dx}{x^2-a^2} = \frac{1}{2a}\ln\Bigl\lvert\frac{x-a}{x+a}\Bigr\rvert + C $ 

\end{enumerate}

\end{document}
